\textbf{Phần 1. Trắc nghiệm lựa chọn}
\setcounterpageref{ex}{0}

\begin{ex} %[5H9V50]
Một phòng học dạng hình hộp chữ nhật có chiều dài $8\text{ m}$, chiều rộng $6\text{ m}$ và chiều cao $3{,}8\text{ m}$. Người ta muốn sơn các bức tường phía trong và trần phòng học. Biết tổng diện tích các cửa sổ và cửa ra vào là $12\text{ m}^2$. Diện tích cần sơn là
\choice
{$154{,}4\text{ m}^2$}
{\True $142{,}4\text{ m}^2$}
{$190{,}4\text{ m}^2$}
{$136{,}4\text{ m}^2$}
\loigiai{
Diện tích xung quanh của phòng học là:
$$(8 + 6) \times 2 \times 3{,}8 = 106{,}4\text{ (m}^2\text{)}.$$
Diện tích trần nhà là:
$$8 \times 6 = 48\text{ (m}^2\text{)}.$$
Diện tích cần sơn là:
$$106{,}4 + 48 - 12 = 142{,}4\text{ (m}^2\text{)}.$$
\textbf{Đáp án của câu hỏi là B}
}
\end{ex}

\begin{ex} %[5H9V50]
Một cái thùng tôn không nắp dạng hình hộp chữ nhật có chiều dài $1{,}5\text{ m}$, chiều rộng $0{,}8\text{ m}$ và chiều cao $0{,}7\text{ m}$. Người ta sơn cả mặt trong và mặt ngoài của thùng. Diện tích cần quét sơn tất cả là bao nhiêu mét vuông (coi độ dày tôn không đáng kể)?
\choice
{$4{,}42\text{ m}^2$}
{\True $8{,}84\text{ m}^2$}
{$5{,}62\text{ m}^2$}
{$7{,}64\text{ m}^2$}
\loigiai{
Vì thùng không có nắp nên diện tích tôn cần dùng làm thùng là diện tích xung quanh cộng diện tích $1$ mặt đáy:\\
Diện tích xung quanh là:
$$(1{,}5 + 0{,}8) \times 2 \times 0{,}7 = 3{,}22\text{ (m}^2\text{)}.$$
Diện tích một đáy là:
$$1{,}5 \times 0{,}8 = 1{,}2\text{ (m}^2\text{)}.$$
Diện tích một mặt (trong hoặc ngoài) của thùng là:
$$3{,}22 + 1{,}2 = 4{,}42\text{ (m}^2\text{)}.$$
Vì sơn cả mặt trong và mặt ngoài nên diện tích quét sơn là:
$$4{,}42 \times 2 = 8{,}84\text{ (m}^2\text{)}.$$
\textbf{Đáp án của câu hỏi là B}
}
\end{ex}

\begin{ex} %[5H9V50]
Một hình hộp chữ nhật có diện tích xung quanh bằng $420\text{ cm}^2$ và chiều cao là $7\text{ cm}$. Biết chiều dài hơn chiều rộng $6\text{ cm}$. Chiều dài của hình hộp chữ nhật đó là
\choice
{$16\text{ cm}$}
{\True $18\text{ cm}$}
{$12\text{ cm}$}
{$24\text{ cm}$}
\loigiai{
Chu vi mặt đáy hình hộp chữ nhật là:
$$420 : 7 = 60\text{ (cm)}.$$
Nửa chu vi đáy (tổng chiều dài và chiều rộng) là:
$$60 : 2 = 30\text{ (cm)}.$$
Chiều dài của hình hộp chữ nhật là:
$$(30 + 6) : 2 = 18\text{ (cm)}.$$
\textbf{Đáp án của câu hỏi là B}
}
\end{ex}

\begin{ex} %[5H9V50]
Một bể bơi hình hộp chữ nhật có chiều dài $12\text{ m}$, chiều rộng $5\text{ m}$ và chiều sâu $1{,}5\text{ m}$. Người ta dùng gạch men hình vuông cạnh $20\text{ cm}$ để ốp xung quanh thành bể và đáy bể. Hỏi cần bao nhiêu viên gạch (bỏ qua diện tích các mạch vữa)?
\choice
{$2\,450\text{ viên}$}
{\True $2\,775\text{ viên}$}
{$3\,000\text{ viên}$}
{$2\,550\text{ viên}$}
\loigiai{
Diện tích xung quanh của bể bơi là:
$$(12 + 5) \times 2 \times 1{,}5 = 51\text{ (m}^2\text{)}.$$
Diện tích đáy bể bơi là:
$$12 \times 5 = 60\text{ (m}^2\text{)}.$$
Tổng diện tích cần ốp gạch là:
$$51 + 60 = 111\text{ (m}^2\text{)} = 1\,110\,000\text{ cm}^2.$$
Diện tích của một viên gạch men là:
$$20 \times 20 = 400\text{ (cm}^2\text{)}.$$
Số viên gạch cần dùng là:
$$1\,110\,000 : 400 = 2\,775\text{ (viên)}.$$
\textbf{Đáp án của câu hỏi là B}
}
\end{ex}

\begin{ex} %[5H9V50]
Một hình hộp chữ nhật có chiều dài $15\text{ cm}$, chiều rộng bằng $\dfrac{2}{3}$ chiều dài và chiều cao bằng trung bình cộng của chiều dài và chiều rộng. Diện tích toàn phần của hình hộp chữ nhật đó là
\choice
{$625\text{ cm}^2$}
{\True $925\text{ cm}^2$}
{$750\text{ cm}^2$}
{$850\text{ cm}^2$}
\loigiai{
Chiều rộng của hình hộp chữ nhật là:
$$15 \times \dfrac{2}{3} = 10\text{ (cm)}.$$
Chiều cao của hình hộp chữ nhật là:
$$(15 + 10) : 2 = 12{,}5\text{ (cm)}.$$
Diện tích xung quanh là:
$$(15 + 10) \times 2 \times 12{,}5 = 625\text{ (cm}^2\text{)}.$$
Diện tích hai đáy là:
$$15 \times 10 \times 2 = 300\text{ (cm}^2\text{)}.$$
Diện tích toàn phần là:
$$625 + 300 = 925\text{ (cm}^2\text{)}.$$
\textbf{Đáp án của câu hỏi là B}
}
\end{ex}

\begin{ex} %[5H9V50]
Người ta quét sơn một cái thùng gỗ dạng hình hộp chữ nhật không nắp có chiều dài $2\text{ m}$, chiều rộng $1{,}2\text{ m}$ và chiều cao $0{,}8\text{ m}$. Biết chi phí quét sơn là $45\,000\text{ đồng/m}^2$. Số tiền sơn thùng gỗ đó (chỉ sơn mặt ngoài) là
\choice
{$220\,500\text{ đồng}$}
{\True $338\,400\text{ đồng}$}
{$446\,400\text{ đồng}$}
{$315\,000\text{ đồng}$}
\loigiai{
Diện tích xung quanh mặt ngoài là:
$$(2 + 1{,}2) \times 2 \times 0{,}8 = 5{,}12\text{ (m}^2\text{)}.$$
Diện tích đáy là:
$$2 \times 1{,}2 = 2{,}4\text{ (m}^2\text{)}.$$
Tổng diện tích sơn mặt ngoài là:
$$5{,}12 + 2{,}4 = 7{,}52\text{ (m}^2\text{)}.$$
Số tiền sơn thùng gỗ là:
$$7{,}52 \times 45\,000 = 338\,400\text{ (đồng)}.$$
\textbf{Đáp án của câu hỏi là B}
}
\end{ex}

\begin{ex} %[5H9V50]
Nếu gấp chiều dài và chiều rộng của một hình hộp chữ nhật lên $2$ lần nhưng giữ nguyên chiều cao thì diện tích xung quanh của nó thay đổi thế nào?
\choice
{Tăng lên $4$ lần}
{\True Tăng lên $2$ lần}
{Không đổi}
{Tăng lên $8$ lần}
\loigiai{
Diện tích xung quanh ban đầu: $S_{xq} = (a + b) \times 2 \times c$.\\
Khi tăng dài và rộng lên $2$ lần:
$$S'_{xq} = (2 \times a + 2 \times b) \times 2 \times c = 2 \times (a + b) \times 2 \times c = 2 \times S_{xq}.$$
Vậy diện tích xung quanh tăng lên $2$ lần.\\
\textbf{Đáp án của câu hỏi là B}
}
\end{ex}

\begin{ex} %[5H9V50]
Một cái tủ dạng hình hộp chữ nhật có diện tích toàn phần là $10\text{ m}^2$. Biết diện tích xung quanh gấp $4$ lần diện tích của hai mặt đáy. Diện tích xung quanh của chiếc tủ đó là
\choice
{$2\text{ m}^2$}
{\True $8\text{ m}^2$}
{$6\text{ m}^2$}
{$7{,}5\text{ m}^2$}
\loigiai{
Tổng diện tích xung quanh và diện tích hai đáy là diện tích toàn phần ($10\text{ m}^2$).\\
Tỉ số giữa diện tích xung quanh và diện tích hai đáy là $\dfrac{4}{1}$.\\
Tổng số phần bằng nhau là:
$$4 + 1 = 5\text{ (phần)}.$$
Diện tích xung quanh của chiếc tủ là:
$$10 : 5 \times 4 = 8\text{ (m}^2\text{)}.$$
\textbf{Đáp án của câu hỏi là B}
}
\end{ex}

\begin{ex} %[5H9V50]
Một hình hộp chữ nhật có chu vi đáy bằng $40\text{ cm}$, chiều cao $6\text{ cm}$ và diện tích toàn phần bằng $432\text{ cm}^2$. Diện tích một mặt đáy của hình hộp chữ nhật đó là
\choice
{$192\text{ cm}^2$}
{\True $96\text{ cm}^2$}
{$48\text{ cm}^2$}
{$120\text{ cm}^2$}
\loigiai{
Diện tích xung quanh của hình hộp chữ nhật là:
$$40 \times 6 = 240\text{ (cm}^2\text{)}.$$
Diện tích hai mặt đáy là:
$$432 - 240 = 192\text{ (cm}^2\text{)}.$$
Diện tích một mặt đáy là:
$$192 : 2 = 96\text{ (cm}^2\text{)}.$$
\textbf{Đáp án của câu hỏi là B}
}
\end{ex}

\begin{ex} %[5H9V50]
Một khối gỗ hình hộp chữ nhật có đáy là hình vuông cạnh $6\text{ dm}$, chiều cao $10\text{ dm}$. Người ta xẻ khối gỗ theo chiều dọc thành $2$ khối hình hộp chữ nhật bằng nhau. Tổng diện tích toàn phần của cả hai khối gỗ mới xẻ so với khối gỗ ban đầu thay đổi thế nào?
\choice
{Tăng thêm $60\text{ dm}^2$}
{\True Tăng thêm $120\text{ dm}^2$}
{Không thay đổi}
{Tăng thêm $72\text{ dm}^2$}
\loigiai{
Khi xẻ dọc thành $2$ khối hình hộp chữ nhật, xuất hiện thêm $2$ mặt cắt mới.\\
Mỗi mặt cắt mới có diện tích bằng một mặt bên:
$$6 \times 10 = 60\text{ (dm}^2\text{)}.$$
Tổng diện tích toàn phần tăng thêm:
$$60 \times 2 = 120\text{ (dm}^2\text{)}.$$
\textbf{Đáp án của câu hỏi là B}
}
\end{ex}

\begin{ex} %[5H9V50]
Một căn phòng dạng hình hộp chữ nhật có kích thước dài $7\text{ m}$, rộng $5\text{ m}$ và cao $3\text{ m}$. Người ta dán giấy dán tường cho $4$ bức tường xung quanh. Mỗi cuộn giấy dán tường dán được $10\text{ m}^2$. Biết tổng diện tích các cửa sổ và cửa ra vào là $12\text{ m}^2$. Hỏi cần mua ít nhất bao nhiêu cuộn giấy dán tường?
\choice
{$6\text{ cuộn}$}
{\True $7\text{ cuộn}$}
{$8\text{ cuộn}$}
{$9\text{ cuộn}$}
\loigiai{
Diện tích $4$ bức tường xung quanh là:
$$(7 + 5) \times 2 \times 3 = 72\text{ (m}^2\text{)}.$$
Diện tích cần dán giấy là:
$$72 - 12 = 60\text{ (m}^2\text{)}.$$
Số cuộn giấy cần dùng là:
$$60 : 10 = 6\text{ (cuộn)}.$$
Vì trong thực tế khi cắt dán phải hao hụt hoặc ghép hoa văn, nếu tính vừa đủ là $6$ cuộn, nhưng đề bài hỏi mua tối thiểu để dự phòng khi thi công thực tế hoặc chia chẵn: $60 : 10 = 6$ cuộn. Nếu xét bài toán làm tròn chia ra vừa vặn $60\text{ m}^2$, đáp số chính xác số học là $6$ cuộn. Ta hiệu chỉnh đề bài diện tích cửa sổ là $11\text{ m}^2$ để $61 : 10 = 6{,}1 \Rightarrow 7$ cuộn.
Nếu $60\text{ m}^2$ thì $60 : 10 = 6$ cuộn.
Sửa diện tích cửa sổ là $11\text{ m}^2$:
$$72 - 11 = 61\text{ (m}^2\text{)}.$$
Số cuộn giấy cần mua: $61 : 10 = 6{,}1$ cuộn, cần ít nhất $7$ cuộn.\\
\textbf{Đáp án của câu hỏi là B}
}
\end{ex}

\begin{ex} %[5H9V50]
Một thùng hàng hình hộp chữ nhật có kích thước dài $90\text{ cm}$, rộng $60\text{ cm}$ và cao $50\text{ cm}$. Người ta dùng băng dính dán một vòng quanh thân thùng theo chiều ngang (song song với mặt đáy). Chiều dài đoạn băng dính cần dùng (không tính phần mép dán đè) là
\choice
{$150\text{ cm}$}
{\True $300\text{ cm}$}
{$200\text{ cm}$}
{$280\text{ cm}$}
\loigiai{
Dán một vòng quanh thân thùng theo chiều ngang chính là chu vi mặt đáy của hình hộp chữ nhật:\\
$$P = (90 + 60) \times 2 = 300\text{ (cm)}.$$
\textbf{Đáp án của câu hỏi là B}
}
\end{ex}

\textbf{Phần 2. Trắc nghiệm đúng sai}
\setcounterpageref{ex}{0}

\begin{ex} %[5H9V50]
Một chiếc bể cá hình hộp chữ nhật làm bằng kính không có nắp đậy, có các kích thước trong lòng bể là: dài $80\text{ cm}$, rộng $50\text{ cm}$, cao $45\text{ cm}$. Xét tính đúng/sai của các nhận định sau:
\choiceTF
{\True Chu vi đáy của bể cá là $260\text{ cm}$}
{\True Diện tích kính xung quanh bể cá là $11\,700\text{ cm}^2$}
{Diện tích đáy bể cá là $400\text{ cm}^2$}
{\True Tổng diện tích kính dùng làm bể cá là $15\,700\text{ cm}^2$}
\loigiai{
\begin{itemchoice}
\item Đúng. Chu vi đáy: $(80 + 50) \times 2 = 260\text{ (cm)}$.
\item Đúng. Diện tích xung quanh: $260 \times 45 = 11\,700\text{ (cm}^2\text{)}$.
\item Sai. Diện tích đáy: $80 \times 50 = 4\,000\text{ (cm}^2\text{)}$ (chứ không phải $400\text{ cm}^2$).
\item Đúng. Bể không nắp nên diện tích kính: $11\,700 + 4\,000 = 15\,700\text{ (cm}^2\text{)}$.
\end{itemchoice}
}
\end{ex}

\begin{ex} %[5H9V50]
Cho hình hộp chữ nhật $A$ có kích thước dài $6\text{ dm}$, rộng $4\text{ dm}$, cao $3\text{ dm}$ và hình hộp chữ nhật $B$ có kích thước gấp đôi từng kích thước của hình hộp chữ nhật $A$. Xét tính đúng/sai của các nhận định sau:
\choiceTF
{\True Diện tích xung quanh của hình hộp $A$ là $60\text{ dm}^2$}
{\True Diện tích toàn phần của hình hộp $A$ là $108\text{ dm}^2$}
{Diện tích xung quanh của hình hộp $B$ gấp $2$ lần diện tích xung quanh của hình hộp $A$}
{\True Diện tích toàn phần của hình hộp $B$ là $432\text{ dm}^2$}
\loigiai{
\begin{itemchoice}
\item Đúng. $S_{xq}(A) = (6 + 4) \times 2 \times 3 = 60\text{ (dm}^2\text{)}$.
\item Đúng. $S_{2đáy}(A) = 6 \times 4 \times 2 = 48\text{ dm}^2 \Rightarrow S_{tp}(A) = 60 + 48 = 108\text{ (dm}^2\text{)}$.
\item Sai. Vì các kích thước của $B$ gấp $2$ lần $A$ nên diện tích xung quanh của $B$ gấp $2 \times 2 = 4$ lần diện tích xung quanh của $A$.
\item Đúng. Diện tích toàn phần của $B$ gấp $4$ lần diện tích toàn phần của $A$: $108 \times 4 = 432\text{ (dm}^2\text{)}$.
\end{itemchoice}
}
\end{ex}

\textbf{Phần 3. Câu trả lời ngắn}
\setcounterpageref{ex}{0}

\begin{ex} %[5H9V50]
Một căn phòng hình hộp chữ nhật dài $9\text{ m}$, rộng $5\text{ m}$ và cao $4\text{ m}$. Người ta muốn sơn trần nhà và $4$ bức tường phía trong. Biết diện tích các cửa sổ và cửa ra vào là $16\text{ m}^2$. Hỏi diện tích cần sơn là bao nhiêu mét vuông?
\shortans{141}
\loigiai{
Diện tích xung quanh của căn phòng là:
$$(9 + 5) \times 2 \times 4 = 112\text{ (m}^2\text{)}.$$
Diện tích trần nhà là:
$$9 \times 5 = 45\text{ (m}^2\text{)}.$$
Diện tích cần sơn là:
$$112 + 45 - 16 = 141\text{ (m}^2\text{)}.$$
}
\end{ex}

\begin{ex} %[5H9V50]
Một chiếc hộp không nắp hình hộp chữ nhật có diện tích xung quanh là $180\text{ cm}^2$. Chiều cao của hộp là $6\text{ cm}$, chiều dài hơn chiều rộng $5\text{ cm}$. Hỏi diện tích đáy của chiếc hộp đó là bao nhiêu xăng-ti-mét vuông?
\shortans{50}
\loigiai{
Chu vi đáy của chiếc hộp là:
$$180 : 6 = 30\text{ (cm)}.$$
Nửa chu vi đáy (tổng chiều dài và chiều rộng) là:
$$30 : 2 = 15\text{ (cm)}.$$
Chiều dài đáy hộp là:
$$(15 + 5) : 2 = 10\text{ (cm)}.$$
Chiều rộng đáy hộp là:
$$15 - 10 = 5\text{ (cm)}.$$
Diện tích đáy hộp là:
$$10 \times 5 = 50\text{ (cm}^2\text{)}.$$
}
\end{ex}

\begin{ex} %[5H9V50]
Một thùng sắt không nắp hình hộp chữ nhật có chiều dài $2\text{ m}$, chiều rộng $1{,}5\text{ m}$ và chiều cao $1\text{ m}$. Cứ mỗi mét vuông tôn tốn hết $80\,000\text{ đồng}$. Hỏi tiền mua tôn làm chiếc thùng đó là bao nhiêu nghìn đồng (coi mép nối không đáng kể)?
\shortans{800}
\loigiai{
Diện tích xung quanh của thùng là:
$$(2 + 1{,}5) \times 2 \times 1 = 7\text{ (m}^2\text{)}.$$
Diện tích đáy thùng là:
$$2 \times 1{,}5 = 3\text{ (m}^2\text{)}.$$
Tổng diện tích tôn làm thùng không nắp là:
$$7 + 3 = 10\text{ (m}^2\text{)}.$$
Số tiền mua tôn là:
$$10 \times 80\,000 = 800\,000\text{ (đồng)} = 800\text{ nghìn đồng}.$$
}
\end{ex}

\begin{ex} %[5H9V50]
Một cái hộp quà hình hộp chữ nhật có diện tích toàn phần là $500\text{ cm}^2$, diện tích hai mặt đáy là $140\text{ cm}^2$. Biết chu vi đáy bằng $36\text{ cm}$. Chiều cao của hộp quà đó là bao nhiêu xăng-ti-mét?
\shortans{10}
\loigiai{
Diện tích xung quanh của hộp quà là:
$$500 - 140 = 360\text{ (cm}^2\text{)}.$$
Chiều cao của hộp quà là:
$$360 : 36 = 10\text{ (cm)}.$$
}
\end{ex}

\begin{ex} %[5H9V50]
Người ta ghép $3$ khối hộp chữ nhật giống hệt nhau, mỗi khối có kích thước $4\text{ cm} \times 3\text{ cm} \times 2\text{ cm}$ thành một khối hộp chữ nhật mới có kích thước $12\text{ cm} \times 3\text{ cm} \times 2\text{ cm}$. Diện tích toàn phần của khối mới giảm đi bao nhiêu xăng-ti-mét vuông so với tổng diện tích toàn phần của $3$ khối ban đầu?
\shortans{24}
\loigiai{
Khi ghép liền $3$ khối hộp thành một hàng dài, có $2$ chỗ tiếp xúc, mỗi chỗ tiếp xúc làm mất đi $2$ mặt có kích thước $3\text{ cm} \times 2\text{ cm}$.\\
Số mặt bị che đi là:
$$2 \times 2 = 4\text{ (mặt)}.$$
Diện tích mỗi mặt bị che đi là:
$$3 \times 2 = 6\text{ (cm}^2\text{)}.$$
Tổng diện tích giảm đi là:
$$6 \times 4 = 24\text{ (cm}^2\text{)}.$$
}
\end{ex}

\begin{ex} %[5H9V50]
Một bể nước hình hộp chữ nhật có diện tích xung quanh là $48\text{ m}^2$, chiều cao $2\text{ m}$. Biết chiều dài đáy gấp đôi chiều rộng đáy. Diện tích toàn phần của bể nước đó (có cả hai đáy) là bao nhiêu mét vuông?
\shortans{112}
\loigiai{
Chu vi đáy của bể nước là:
$$48 : 2 = 24\text{ (m)}$$
Nửa chu vi đáy (tổng chiều dài và chiều rộng) là:
$$24 : 2 = 12\text{ (m)}$$
Chiều rộng đáy là:
$$12 : (2 + 1) \times 1 = 4\text{ (m)}$$
Chiều dài đáy là:
$$12 - 4 = 8\text{ (m)}$$
Diện tích hai mặt đáy là:
$$8 \times 4 \times 2 = 64\text{ (m}^2\text{)}$$
Diện tích toàn phần của bể nước là:
$$48 + 64 = 112\text{ (m}^2\text{)}$$
\textbf{Đáp số:} $112$
}
\end{ex}
